\chapter{Euclid Revisited}\label{ch:euclid}

\begin{center}\itshape
There is no royal road to geometry.
\upshape(Euclid to King Ptolemy, as reported by Proclus in his \emph{Commentary on the First Book of Euclid's Elements})
\end{center}

Around 1899 Frank Morley, an English mathematician teaching at Haverford College in Pennsylvania, was studying curves
tangent to the three sides of a triangle. Along the way he noticed something that had nothing to do with curves. Take any
triangle and trisect each of its angles. Near each side two trisectors meet, one from each end of that side. The three
meeting points form an \emph{equilateral} triangle, whatever the shape of the original one. Morley mentioned it to friends
and moved on. The result spread by word of mouth, and when it finally appeared in print in the 1900s it caused
astonishment. Euclid had all the tools to find it, and for two thousand years nobody had. Since then dozens of proofs
have been published, by trigonometry, by pure geometry, by complex numbers, and in 1998 Alain Connes gave one that sees it
as a fact about the group of affine maps of a line.

Morley's theorem is the most famous of the late surprises of triangle geometry, but not the only one. In 1765 Euler proved that
three of the classical centres of a triangle always lie on one line. In 1822 Karl Feuerbach found that the circle through the midpoints of the sides
of a triangle touches its incircle and its three excircles; Olry Terquem showed in 1842 that it passes through nine special points. Legend credits Napoleon Bonaparte with a theorem about
equilateral triangles built on the sides of any triangle, although nobody has ever found a source in which he states it.

This chapter proves all of these, and the tool is the one from Tome~0, Chapter~2: complex numbers. A point of the plane is a
complex number, a rotation is a multiplication, and a classical theorem becomes a short identity. The art is in choosing
coordinates well, and the best choice for a triangle is usually to put its circumscribed circle on the unit circle, where
conjugation is the same as inversion: $\bar z=1/z$.

\begin{think}
In a triangle, the \emph{centroid} $G$ is where the three medians meet, the \emph{circumcentre} $O$ is the centre of the circle
through the vertices, and the \emph{orthocentre} $H$ is where the three altitudes meet. Draw a few triangles carefully. Do the
three points have anything in common? If they are on one line, guess the ratio $OG:GH$.
\end{think}

\section{The plane as the complex numbers}

We identify the Euclidean plane with $\C$: the point $(x,y)$ is $z=x+iy$, the distance from $z$ to $w$ is $|z-w|$, and
multiplication by $e^{i\theta}$ is the rotation by $\theta$ about $0$. Three tests do most of the work.

\begin{proposition}[three tests]\label{prop:tests}
Let $a,b,z\in\C$ with $a\ne b$.
\begin{enumerate}
\item $z$ lies on the line $ab$ if and only if $\dfrac{z-a}{b-a}$ is real.
\item The line $az$ (for $z\ne a$) is perpendicular to the line $ab$ if and only if $\dfrac{z-a}{b-a}$ is purely imaginary.
\item The signed area of the triangle $abc$, positive when $a,b,c$ go round counterclockwise, is
\[
[abc]=\tfrac12\operatorname{Im}\bigl(\overline{(b-a)}(c-a)\bigr)=\frac i4\begin{vmatrix}a&\bar a&1\\ b&\bar b&1\\ c&\bar c&1\end{vmatrix}.
\]
\end{enumerate}
\end{proposition}
\begin{proof}
The argument of $(z-a)/(b-a)$ is the angle from the vector $b-a$ to the vector $z-a$. It is $0$ or $\pi$ exactly when the
vectors are parallel, which gives (1), and $\pm\pi/2$ exactly when they are perpendicular, which gives (2). For (3), write
$u=b-a=u_1+iu_2$ and $v=c-a=v_1+iv_2$. The signed area is $\frac12(u_1v_2-u_2v_1)$, the determinant of S1, Chapter~2, and
$\bar uv=(u_1v_1+u_2v_2)+i(u_1v_2-u_2v_1)$. For the determinant, subtract the first row from the others; the expansion
along the last column leaves $(b-a)\overline{(c-a)}-\overline{(b-a)}(c-a)=-2i\operatorname{Im}(\bar uv)$, and multiplying by
$i/4$ gives $\frac12\operatorname{Im}(\bar uv)$.
\end{proof}

When $a$ and $b$ lie on the unit circle, $\bar a=1/a$ and $\bar b=1/b$, and the tests become linear equations.

\begin{lemma}[chords of the unit circle]\label{lem:chord}
Let $a\ne b$ with $|a|=|b|=1$.
\begin{enumerate}
\item A point $z$ lies on the line $ab$ if and only if $z+ab\,\bar z=a+b$.
\item The foot of the perpendicular from any point $z$ to the line $ab$ is $f=\frac12\bigl(a+b+z-ab\,\bar z\bigr)$, and the mirror
image of $z$ in the line is $2f-z=a+b-ab\,\bar z$.
\end{enumerate}
\end{lemma}
\begin{proof}
(1) By Proposition~\ref{prop:tests}, $z$ is on the line if and only if $\frac{z-a}{b-a}$ equals its conjugate
$\frac{\bar z-1/a}{1/b-1/a}=\frac{ab\bar z-b}{a-b}$, that is $z-a=b-ab\bar z$.

(2) The conjugate of $f$ is $\frac12(1/a+1/b+\bar z-z/(ab))$, so $ab\bar f=\frac12(b+a+ab\bar z-z)$ and $f+ab\bar f=a+b$: the point $f$
is on the line. Next, $w=a+b-z-ab\bar z$ satisfies $\bar w=w/(ab)$, and $\overline{b-a}=(a-b)/(ab)$, so the conjugate of
$w/(b-a)$ is $w/(a-b)=-w/(b-a)$. Thus $(f-z)/(b-a)=w/(2(b-a))$ is purely imaginary and $zf$ is perpendicular to $ab$.
\end{proof}

\section{The centres of a triangle and Euler's line}

From now on $ABC$ is a triangle, and we scale and move it so that its circumscribed circle is the unit circle: the circumcentre is
$O=0$, the circumradius is $R=1$, and the vertices are complex numbers $a,b,c$ of modulus $1$. The centroid is
$g=\frac13(a+b+c)$, since the medians meet two thirds of the way from each vertex to the midpoint of the opposite side (indeed
$g=\frac13a+\frac23\cdot\frac{b+c}2$, and similarly for $b$ and $c$).

\begin{proposition}[the orthocentre]\label{prop:orthocentre}
The three altitudes of $ABC$ meet at the point $h=a+b+c$.
\end{proposition}
\begin{proof}
We show that $h$ lies on the altitude from $A$, the line through $a$ perpendicular to $bc$; the other two follow by symmetry. Now
$h-a=b+c$, and the conjugate of $\frac{b+c}{b-c}$ is $\frac{1/b+1/c}{1/b-1/c}=\frac{c+b}{c-b}=-\frac{b+c}{b-c}$. So $\frac{h-a}{b-c}$
is purely imaginary (or zero, if $h=a$), and by Proposition~\ref{prop:tests} the point $h$ is on the altitude.
\end{proof}

\begin{theorem}[Euler, 1765]\label{thm:eulerline}
In every triangle, the circumcentre $O$, the centroid $G$ and the orthocentre $H$ lie on one line, with $G$ between $O$ and $H$ and
$GH=2\,OG$. They coincide exactly when the triangle is equilateral.
\end{theorem}
\begin{proof}
With $O=0$ we have $h=a+b+c=3g$: the vectors from $O$ to $G$ and to $H$ point the same way, and $OH=3\,OG$. If $g=0$, all three
points coincide; then the median from $A$ passes through $O$, so it is perpendicular to $BC$ (the perpendicular bisector of a chord
passes through the centre), hence $AB=AC$, and similarly $BA=BC$. Conversely, an equilateral triangle has all its centres at one point.
\end{proof}

The line through $O$, $G$ and $H$ is the \emph{Euler line}. The midpoint of $OH$ is the point $n=\frac12(a+b+c)$, and it is the centre
of a remarkable circle.

\begin{theorem}[the nine-point circle]\label{thm:ninepoint}
In every triangle, the following nine points lie on one circle, of centre the midpoint $N$ of $OH$ and of radius $R/2$: the midpoints of
the three sides, the feet of the three altitudes, and the midpoints of the segments joining the orthocentre to the vertices.
\end{theorem}
\begin{proof}
We measure distances from $n=\frac12(a+b+c)$. The midpoint of $BC$ is $\frac12(b+c)=n-\frac a2$, at distance $\frac12|a|=\frac12$.
The midpoint of $AH$ is $\frac12(a+h)=n+\frac a2$, again at distance $\frac12$. By Lemma~\ref{lem:chord} with $z=a$, the foot of the
altitude from $A$ is
\[
d=\tfrac12\bigl(b+c+a-bc\,\bar a\bigr)=n-\frac{bc}{2a},
\]
at distance $\frac12|bc/a|=\frac12$. The other six points follow by symmetry.
\end{proof}

Figure~\ref{fig:euler} shows the nine points. The midpoints of the sides and of the segments $AH$, $BH$, $CH$ come in diametrically
opposite pairs ($n\pm a/2$), so the segment from the midpoint of $BC$ to the midpoint of $AH$ is a diameter of the nine-point circle.

\begin{figure}[ht]
\centering
\fig{ch01-euler}
\caption{The Euler line (red) through $O$, $G$, $N$ and $H$, with $OG:GH=1:2$ and $N$ the midpoint of $OH$, and the nine-point circle
(green) of radius $R/2$ through the midpoints of the sides, the feet of the altitudes (dashed) and the midpoints of $AH$, $BH$, $CH$.
The triangle has angles $75^\circ$, $60^\circ$ and $45^\circ$.}
\label{fig:euler}
\end{figure}

\begin{remark}
In 1822 Karl Wilhelm Feuerbach proved that the nine-point circle touches the inscribed circle and the three escribed circles. We prove
this in Chapter~3.
\end{remark}

\subsection*{The incentre}

The incentre $I$, where the three angle bisectors meet, is less symmetric in $a,b,c$: it needs square roots.

\begin{proposition}[the incentre]\label{prop:incentre}
One can choose square roots $u,v,w$ of $a,b,c$ such that the midpoints of the arcs $BC$, $CA$, $AB$ not containing the opposite vertex are
$-vw$, $-wu$, $-uv$. Then the incentre is $i=-(uv+vw+wu)$.
\end{proposition}
\begin{proof}
Rotate so that $a=e^{2i\alpha}$, $b=e^{2i\beta}$, $c=e^{2i\gamma}$ with $\alpha<\beta<\gamma<\alpha+\pi$, the vertices in counterclockwise
order. The arc from $A$ to $B$ counterclockwise does not contain $C$, and its midpoint is $e^{i(\alpha+\beta)}$; similarly the
midpoint of the arc from $B$ to $C$ is $e^{i(\beta+\gamma)}$, and the arc from $C$ counterclockwise to $A$, that is from angle $2\gamma$
to $2\alpha+2\pi$, has midpoint $e^{i(\gamma+\alpha+\pi)}=-e^{i(\gamma+\alpha)}$. Put $u=e^{i\alpha}$, $v=-e^{i\beta}$, $w=e^{i\gamma}$; the three
midpoints become $-uv$, $-vw$, $-wu$, as claimed.

By the inscribed angle theorem, equal arcs subtend equal angles, so the bisector of the angle at $A$ passes through the midpoint $m_a=-vw$
of the arc $BC$ not containing $A$. It remains to check that $i=-(uv+vw+wu)$ is on the line through $u^2$ and $-vw$, that is, by
Lemma~\ref{lem:chord}, that $i-u^2vw\,\bar i=u^2-vw$. Since $\bar u=1/u$ and so on, $\bar i=-\bigl(\frac1{uv}+\frac1{vw}+\frac1{wu}\bigr)
=-\frac{u+v+w}{uvw}$, hence
\[
i-u^2vw\,\bar i=-(uv+vw+wu)+u(u+v+w)=u^2-vw .
\]
By symmetry $i$ is on all three bisectors, so it is the incentre.
\end{proof}

\section{Ceva's theorem}

A \emph{cevian} is a line through a vertex of a triangle and a point of the opposite side line, named after Giovanni Ceva, who
published the following criterion in 1678 (it was already known to Yusuf al-Mu'taman ibn H\=ud, an eleventh-century king of Zaragoza
and a fine mathematician; see Hogendijk in the sources). We
use signed ratios: for $D$ on the line $BC$, $\frac{BD}{DC}$ is the real number $t$ with $D-B=t(C-D)$; it is positive exactly when $D$ is
between $B$ and $C$.

\begin{theorem}[Ceva]\label{thm:ceva}
Let $D,E,F$ be points of the lines $BC$, $CA$, $AB$, different from the vertices. The lines $AD$, $BE$, $CF$ meet in one point or are
parallel if and only if
\[
\frac{BD}{DC}\cdot\frac{CE}{EA}\cdot\frac{AF}{FB}=1 .
\]
\end{theorem}
\begin{proof}
Every point $P$ of the plane can be written $P=\frac{\alpha A+\beta B+\gamma C}{\alpha+\beta+\gamma}$ with real weights, unique up to a
common factor (its \emph{barycentric coordinates}): the vectors $B-A$ and $C-A$ form a basis, and $P-A=y(B-A)+z(C-A)$ gives
$(\alpha,\beta,\gamma)=(1-y-z,y,z)$.

Suppose the cevians meet at $P$ with weights $(\alpha,\beta,\gamma)$. If $\beta+\gamma=0$, then $\alpha\ne0$ and $P=A+\frac\beta\alpha(B-C)$.
Here $\beta\ne0$, since otherwise $P=A$ and the line $BE$ would pass through $A$, forcing $E=A$. So the line $AP$ would be parallel to
$BC$; but it contains the point $D$ of $BC$, so it would be $BC$ itself, which does not contain $A$. Hence $\beta+\gamma\neq0$, and the point $D'=\frac{\beta B+\gamma C}{\beta+\gamma}$ is on $BC$ and on the line $AP$, because
$P=\frac{\alpha A+(\beta+\gamma)D'}{\alpha+\beta+\gamma}$. Hence $D'=D$, and $D-B=\frac{\gamma}{\beta+\gamma}(C-B)$, $C-D=\frac{\beta}{\beta+\gamma}(C-B)$ give
$\frac{BD}{DC}=\frac\gamma\beta$. Likewise $\frac{CE}{EA}=\frac\alpha\gamma$ and $\frac{AF}{FB}=\frac\beta\alpha$, and the product is $1$. The case of
three parallel cevians is Exercise~\ref{ex:cevaparallel}.

Conversely, assume the product is $1$. Choose $\alpha=1$, $\beta=\frac{AF}{FB}$ and $\gamma=\frac{EA}{CE}$, so that
$\frac{AF}{FB}=\frac\beta\alpha$ and $\frac{CE}{EA}=\frac\alpha\gamma$; the hypothesis gives $\frac{BD}{DC}=\frac\gamma\beta$. None of $\alpha+\beta$,
$\beta+\gamma$, $\gamma+\alpha$ vanishes, since a ratio equal to $-1$ would put the point at infinity. If $\alpha+\beta+\gamma\ne0$,
the point $P=\frac{\alpha A+\beta B+\gamma C}{\alpha+\beta+\gamma}$ is on the line joining $A$ to $D'=\frac{\beta B+\gamma C}{\beta+\gamma}$, as in the
first part; and $D'$ is the point of $BC$ with $\frac{BD'}{D'C}=\frac\gamma\beta=\frac{BD}{DC}$, so $D'=D$ (the signed ratio determines the point).
In the same way $P$ is on $BE$ and $CF$.
If $\alpha+\beta+\gamma=0$, then $D-A=\frac{\beta B+\gamma C}{\beta+\gamma}-A=\frac{\beta B+\gamma C+\alpha A}{\beta+\gamma}$, and the same
vector $\alpha A+\beta B+\gamma C$ (which does not depend on the origin when the weights sum to $0$) is parallel to all three
cevians.
\end{proof}

\begin{example}\label{exa:gergonne}
The medians meet, since all three ratios are $1$. More interesting: let the inscribed circle touch $BC$, $CA$, $AB$ at $D,E,F$. The two
tangents from a point to a circle have equal length, so $AE=AF=x$, $BF=BD=y$, $CD=CE=z$, and the product in Ceva's theorem is
$\frac yz\cdot\frac zx\cdot\frac xy=1$. The lines $AD$, $BE$, $CF$ meet in a point, the \emph{Gergonne point}.
\end{example}

\section{Napoleon's theorem}

The cube roots of unity $1$, $\omega=e^{2\pi i/3}$, $\omega^2$ satisfy $1+\omega+\omega^2=0$, and they recognise equilateral
triangles.

\begin{lemma}\label{lem:equilateral}
The points $a,b,c$ form an equilateral triangle traversed counterclockwise if and only if $a+\omega b+\omega^2c=0$.
\end{lemma}
\begin{proof}
The triangle $abc$ is equilateral and counterclockwise exactly when $c-a$ is $b-a$ turned by $+60^\circ$, that is
$c-a=e^{i\pi/3}(b-a)$. Since $e^{i\pi/3}=-\omega^2$, this reads $c=a+\omega^2a-\omega^2b=-\omega a-\omega^2b$, that is
$c+\omega a+\omega^2b=0$. Multiplying by $\omega^2$ and using $\omega^3=1$ gives $a+\omega b+\omega^2c=0$; all steps reverse.
\end{proof}

\begin{think}
On each side of a triangle, build an equilateral triangle pointing outwards, and mark its centre. What shape do the three centres form?
Try it with a very flat triangle.
\end{think}

\begin{theorem}[Napoleon's theorem]\label{thm:napoleon}
The centres of the three equilateral triangles erected outwards on the sides of any triangle form an equilateral triangle, whose centre is
the centroid of the original triangle.
\end{theorem}
\begin{proof}
Label the triangle counterclockwise. The outward triangle on $BC$ lies on the other side of $BC$ from $A$, so its vertices $c$, $b$, $p$
go round counterclockwise, and by Lemma~\ref{lem:equilateral}, $c+\omega b+\omega^2p=0$, that is $p=-\omega c-\omega^2b$ (multiply by
$\omega$). Its centre is the average of its vertices:
\[
3n_a=b+c+p=(1-\omega^2)b+(1-\omega)c.
\]
In the same way the outward triangles on $CA$ and $AB$ have centres with $3n_b=(1-\omega^2)c+(1-\omega)a$ and
$3n_c=(1-\omega^2)a+(1-\omega)b$. In $3(n_a+\omega n_b+\omega^2n_c)$ the coefficient of $a$ is $\omega(1-\omega)+\omega^2(1-\omega^2)=
\omega-\omega^4=0$, that of $b$ is $(1-\omega^2)+\omega^2(1-\omega)=1-\omega^3=0$, and that of $c$ is $(1-\omega)+\omega(1-\omega^2)=0$.
By Lemma~\ref{lem:equilateral}, $n_an_bn_c$ is equilateral. Finally $n_a+n_b+n_c=\frac13(2-\omega-\omega^2)(a+b+c)=a+b+c$, so the centre of
the Napoleon triangle is $\frac13(a+b+c)$.
\end{proof}

\begin{figure}[ht]
\centering
\fig{ch01-napoleon}
\caption{Napoleon's theorem: the centres of the three outward equilateral triangles (blue) form an equilateral triangle (red), centred at
the centroid of $ABC$.}
\label{fig:napoleon}
\end{figure}

The theorem appeared in print in 1825, in a problem of the \emph{Ladies' Diary}, four years after Napoleon's death; the attribution came
much later (see Wetzel's article in the sources). Napoleon was a keen amateur of geometry, and Coxeter and Greitzer retell a story
in which, after he explained some geometry to Lagrange and Laplace, Laplace replied: ``The last thing we want from you, General, is a
lesson in geometry.'' None of this connects him with the theorem.

\section{A gem: Morley's theorem}

\begin{think}
In Figure~\ref{fig:morley} the angles of $ABC$ are trisected. Before reading on, measure the three sides of the small red triangle on
your own drawing of some other triangle.
\end{think}

\begin{theorem}[Morley, 1899]\label{thm:morley}
In any triangle, the three points where adjacent angle trisectors meet are the vertices of an equilateral triangle, of side
\[
8R\sin\frac A3\sin\frac B3\sin\frac C3 ,
\]
where $R$ is the circumradius and $A$, $B$, $C$ are the angles.
\end{theorem}

The proof uses two facts of trigonometry. The \emph{law of sines}: in a triangle with circumradius $R$, the side opposite an angle
$\theta$ has length $2R\sin\theta$ (the side is a chord of the circumcircle seen under the inscribed angle $\theta$). And the identity
\begin{equation}\label{eq:sin3}
\sin3x=4\sin x\,\sin(60^\circ-x)\,\sin(60^\circ+x),
\end{equation}
which holds because $\sin(60^\circ-x)\sin(60^\circ+x)=\sin^260^\circ-\sin^2x=\frac34-\sin^2x$ and $\sin3x=3\sin x-4\sin^3x$.

\begin{proof}
Write $A=3\alpha$, $B=3\beta$, $C=3\gamma$, so that $\alpha+\beta+\gamma=60^\circ$. Let $P$ be the meeting point of the trisectors from $B$ and
$C$ adjacent to $BC$, and $R'$ that of the trisectors from $A$ and $B$ adjacent to $AB$ (Figure~\ref{fig:morley}, where $R'$ is labelled
$R$).

In the triangle $BPC$ the angles at $B$ and $C$ are $\beta$ and $\gamma$, so the angle at $P$ is $180^\circ-\beta-\gamma$, whose sine is
$\sin(\beta+\gamma)=\sin(60^\circ-\alpha)$. The law of sines in $BPC$ and in $ABC$ gives
\[
BP=\frac{BC\,\sin\gamma}{\sin(60^\circ-\alpha)}=\frac{2R\sin3\alpha\,\sin\gamma}{\sin(60^\circ-\alpha)}=8R\sin\alpha\sin\gamma\sin(60^\circ+\alpha),
\]
by \eqref{eq:sin3}. In the triangle $AR'B$ the angles at $A$ and $B$ are $\alpha$ and $\beta$, and the same computation gives
\[
BR'=\frac{AB\,\sin\alpha}{\sin(60^\circ-\gamma)}=\frac{2R\sin3\gamma\,\sin\alpha}{\sin(60^\circ-\gamma)}=8R\sin\alpha\sin\gamma\sin(60^\circ+\gamma).
\]
Put $k=8R\sin\alpha\sin\gamma$. The rays $BP$ and $BR'$ are the two outer trisectors of the angle at $B$, so the angle $PBR'$ is $\beta$.

Now consider an auxiliary triangle $XYZ$ with angles $\beta$ at $X$, $60^\circ+\gamma$ at $Y$ and $60^\circ+\alpha$ at $Z$; these add up
to $180^\circ$, so it exists, and we scale it to have circumdiameter $k$. By the law of sines its sides are $XZ=k\sin(60^\circ+\gamma)$,
$XY=k\sin(60^\circ+\alpha)$ and $YZ=k\sin\beta$. So the triangles $BPR'$ and $XYZ$ have two sides and the included angle equal
($BP=XY$, $BR'=XZ$, angle $\beta$), hence they are congruent, and
\[
PR'=YZ=k\sin\beta=8R\sin\alpha\sin\beta\sin\gamma .
\]
This expression is symmetric in $\alpha,\beta,\gamma$. The same argument at the vertices $C$ and $A$ shows that the other two sides of the
Morley triangle have the same length, which proves the theorem.
\end{proof}

\begin{figure}[ht]
\centering
\fig{ch01-morley}
\caption{Morley's theorem. The adjacent trisectors (blue) meet at $P$, $Q$, $R$, which form an equilateral triangle. The side is
$8R\sin\frac A3\sin\frac B3\sin\frac C3$, checked in Python to $10^{-15}$ for this triangle.}
\label{fig:morley}
\end{figure}

The proof also tells us the angles of the small triangles around the Morley triangle: in $BPR'$ they are $\beta$, $60^\circ+\gamma$ and
$60^\circ+\alpha$. John Conway turned this into a famous ``backwards'' proof: start from an equilateral triangle, glue onto its sides
triangles with exactly these angles, check that the pieces fit together into a triangle with angles $3\alpha$, $3\beta$, $3\gamma$, and
conclude because a triangle is determined up to similarity by its angles. In 1998 Alain Connes found a proof of a different kind. Let
$f,g,h$ be the rotations about $A$, $B$, $C$ by $\frac23A$, $\frac23B$, $\frac23C$. Their cubes are rotations by twice the angles of the
triangle, and the product of these is the identity (each is a product of two reflections in sides of the triangle, which cancel).
From $f^3g^3h^3=1$ alone, a computation in the group of maps $z\mapsto az+b$ shows that the fixed points of $fg$, $gh$, $hf$, which are the
Morley points, form an equilateral triangle. It works over any field in which the relevant cube roots exist.

\begin{beauty}
A triangle has no preferred direction and no preferred size, and nothing about its angles is special. Yet five centuries of geometers
kept finding that it carries rigid, perfectly symmetric structures: a line through three centres, a circle through nine points, an
equilateral triangle among its trisectors. Complex numbers explain why: the symmetric functions $a+b+c$ and $ab+bc+ca$ of the vertices
are the hidden coordinates of the triangle, and every theorem of this chapter is an identity about them.
\end{beauty}

\section{Exercises}

\begin{exercise}[$\star$]\label{ex:reflectH}
Show that the mirror images of the orthocentre in the three sides of a triangle lie on its circumcircle.
\end{exercise}

\begin{exercise}[$\star\star$]\label{ex:simson}
Let $P$ be a point of the circumcircle of $ABC$. Show that the feet of the perpendiculars from $P$ to the three side lines are collinear
(the \emph{Simson line} of $P$).
\end{exercise}

\begin{exercise}[$\star\star$]\label{ex:eulerOI}
Prove Euler's formula $OI^2=R(R-2r)$, where $r$ is the inradius, and deduce that $R\ge2r$, with equality only for the equilateral triangle.
\end{exercise}

\begin{exercise}[$\star$]\label{ex:viviani}
Show that the sum of the distances from a point inside an equilateral triangle to its three sides does not depend on the point.
\end{exercise}

\begin{exercise}[$\star\star$]\label{ex:fermatpoint}
Let $ABC$ be a triangle with all angles less than $120^\circ$. Show that the point $P$ minimising $PA+PB+PC$ is the point from which the
three sides are seen under angles of $120^\circ$. (Rotate the triangle $ABP$ by $60^\circ$ about $B$.)
\end{exercise}

\begin{exercise}[$\star\star$]\label{ex:vanaubel}
On the four sides of any quadrilateral $ABCD$, build squares outwards. Show that the segments joining the centres of opposite squares
are equal in length and perpendicular (van Aubel's theorem).
\end{exercise}

\begin{exercise}[$\star\star$]\label{ex:miquel}
Points $D$, $E$, $F$ are chosen on the sides $BC$, $CA$, $AB$ of a triangle. Show that the circles $AEF$, $BFD$ and $CDE$ pass through
one point (Miquel's theorem).
\end{exercise}

\begin{exercise}[$\star\star\star$]\label{ex:steinerlehmus}
Show that a triangle with two equal internal angle bisectors (measured from the vertex to the opposite side) is isosceles (the
Steiner--Lehmus theorem).
\end{exercise}

\begin{exercise}[$\star\star$]\label{ex:gausslucas}
Show that the roots of the derivative of a non-constant complex polynomial lie in the convex hull of the roots of the polynomial (the
Gauss--Lucas theorem).
\end{exercise}

\begin{exercise}[$\star\star$]\label{ex:routh}
In a triangle of area $1$, each side is divided in the ratio $1:2$ (going round the triangle in the same direction), and each vertex is
joined to the division point of the opposite side. Show that the three cevians enclose a triangle of area $\frac17$.
\end{exercise}

\begin{exercise}[$\star$]\label{ex:latticeequilateral}
Show that the area of a triangle whose vertices have integer coordinates is a multiple of $\frac12$, and deduce that no equilateral
triangle has its three vertices at points with integer coordinates.
\end{exercise}

\begin{exercise}[$\star$]\label{ex:cevaparallel}
Finish the proof of Ceva's theorem: if the three cevians $AD$, $BE$, $CF$ are parallel, show that the product of the ratios is $1$.
\end{exercise}

\section{Hints}

\begin{hint}{ex:reflectH}
Lemma~\ref{lem:chord} gives the mirror image of $h=a+b+c$ in the line $bc$.
\end{hint}
\begin{hint}{ex:simson}
Compute the feet with Lemma~\ref{lem:chord} and factor the differences of two of them.
\end{hint}
\begin{hint}{ex:eulerOI}
Let the bisector $AI$ meet the circumcircle again at $M$. Show that $MB=MI$, and compute the power of $I$ with respect to the circumcircle.
\end{hint}
\begin{hint}{ex:viviani}
Cut the triangle into three triangles with apex at the point.
\end{hint}
\begin{hint}{ex:fermatpoint}
After the rotation, $PA+PB+PC$ is the length of a broken line from a fixed point to $C$.
\end{hint}
\begin{hint}{ex:vanaubel}
For a counterclockwise polygon, the centre of the outer square on the side from $z$ to $w$ is $\frac{z+w}2-i\frac{w-z}2$.
\end{hint}
\begin{hint}{ex:miquel}
Let the circles $AEF$ and $BFD$ meet again at $M$, and compute the angle $DME$ using cyclic quadrilaterals.
\end{hint}
\begin{hint}{ex:steinerlehmus}
Compute the length of the bisector with Stewart's theorem, then suppose $b<c$ and compare.
\end{hint}
\begin{hint}{ex:gausslucas}
Write $p'/p$ as a sum over the roots.
\end{hint}
\begin{hint}{ex:routh}
An affine map preserves ratios of areas, so you may take $A=(0,0)$, $B=(1,0)$, $C=(0,1)$.
\end{hint}
\begin{hint}{ex:latticeequilateral}
Proposition~\ref{prop:tests}(3), and the area of an equilateral triangle of side $s$.
\end{hint}
\begin{hint}{ex:cevaparallel}
Choose coordinates in which the three cevians are vertical, and read the ratios on the $x$-axis.
\end{hint}

\section{Solutions}

\begin{solution}{ex:reflectH}
\idea A computation on the unit circle, which Lemma~\ref{lem:chord} makes one line long.
\writeup With $O=0$ and $R=1$, the mirror image of $h=a+b+c$ in the line $bc$ is $b+c-bc\,\bar h=b+c-bc\bigl(\frac1a+\frac1b+\frac1c\bigr)=-\frac{bc}a$,
which has modulus $1$. The same holds for the other two sides.
\end{solution}

\begin{solution}{ex:simson}
\idea Three points are collinear when the ratio of two difference vectors is real, and on the unit circle conjugates are easy.
\writeup Put the circumcircle on the unit circle; $|p|=1$. By Lemma~\ref{lem:chord}, the feet on $BC$, $CA$, $AB$ are
$d_a=\frac12(p+b+c-bc/p)$, $d_b=\frac12(p+c+a-ca/p)$ and $d_c=\frac12(p+a+b-ab/p)$. Then
\[
d_a-d_c=\tfrac12(c-a)\bigl(1-\tfrac bp\bigr)=\frac{(c-a)(p-b)}{2p},\qquad d_b-d_c=\frac{(c-b)(p-a)}{2p},
\]
and the ratio is $\lambda=\frac{(c-a)(p-b)}{(c-b)(p-a)}$. Using $\bar z=1/z$, $\overline{c-a}=\frac{a-c}{ac}$ and so on, we get
$\bar\lambda=\frac{(a-c)(b-p)\,bc\cdot ap}{ac\cdot bp\,(b-c)(a-p)}=\lambda$. So $\lambda$ is real and the three feet are collinear. (If two feet coincide
there is nothing to prove.) The converse also holds: the feet are collinear only when $P$ is on the circumcircle.
\end{solution}

\begin{solution}{ex:eulerOI}
\idea The power of a point turns a distance to the centre into a product of two lengths along a chord, and the bisector from $A$ is a
chord through $I$.
\writeup Let $M$ be the second intersection of $AI$ with the circumcircle, the midpoint of the arc $BC$ (Proposition~\ref{prop:incentre}).
The angle $MBC$ equals the angle $MAC=A/2$ (same arc $MC$), so $\angle MBI=\frac A2+\frac B2$. The angle $BIM$ is an exterior angle of the
triangle $ABI$, so $\angle BIM=\frac A2+\frac B2$. The triangle $MBI$ is isosceles: $MI=MB$. The chord $MB$ is seen from $A$ under the
angle $A/2$, so $MB=2R\sin\frac A2$. Also $AI=r/\sin\frac A2$, from the right triangle formed by $A$, $I$ and the point where the incircle
touches $AB$. The chord $AM$ passes through $I$, which is inside the circle, so the product $IA\cdot IM$ equals $R^2-OI^2$: the points of the line
$AM$ are $I+tu$ with $u$ a unit vector, they lie on the circumcircle (centre $O$, radius $R$) exactly when $t^2+2t\langle u,I-O\rangle+OI^2-R^2=0$,
the two roots are $\pm IA$ and $\mp IM$ (the points $A$ and $M$ lie on opposite sides of $I$), and their product is $OI^2-R^2$. Hence
\[
R^2-OI^2=\frac r{\sin(A/2)}\cdot2R\sin\frac A2=2Rr,
\]
that is $OI^2=R(R-2r)$. As $OI^2\ge0$, $R\ge2r$, with equality exactly when $I=O$; then the bisectors are perpendicular bisectors and
the triangle is equilateral.
\end{solution}

\begin{solution}{ex:viviani}
\idea Area is additive, and each piece has the side of the triangle as base.
\writeup Let the side be $s$ and the distances be $x,y,z$. The three triangles with apex $P$ have areas $\frac12sx$, $\frac12sy$, $\frac12sz$,
which add up to the area $\frac12sh$ of the whole triangle, where $h$ is its height. So $x+y+z=h$.
\end{solution}

\begin{solution}{ex:fermatpoint}
\idea A sum of three lengths is the length of a broken line once the pieces are put end to end, and a rotation by $60^\circ$ about $B$ does
that, because it turns $BP$ into a side of an equilateral triangle.
\writeup Let $\rho$ be the rotation by $60^\circ$ about $B$ taking $A$ to a point $A'$ on the other side of $AB$ from $C$, and let
$P'=\rho(P)$. The triangle $BPP'$ is equilateral, so $PP'=PB$, and $P'A'=PA$. Thus $PA+PB+PC=A'P'+P'P+PC\ge A'C$, a length that does not
depend on $P$. Equality holds when $A'$, $P'$, $P$, $C$ lie in this order on a segment. Then $\angle BPC=180^\circ-\angle BPP'=120^\circ$
and $\angle APB=\angle A'P'B=180^\circ-\angle BP'P=120^\circ$; the third angle is $120^\circ$ too. Conversely, the point $P$ seeing the sides
under $120^\circ$ exists when all angles are less than $120^\circ$ (it is the second intersection of the circles on which $AB$ and $BC$ are
seen under $120^\circ$), and for it the four points are aligned, so it achieves the minimum $A'C$.
\end{solution}

\begin{solution}{ex:vanaubel}
\idea Write the four centres with complex numbers and look for a multiplication by $\pm i$.
\writeup Take $ABCD$ counterclockwise; the outside of the side from $z$ to $w$ is on the right, the direction $-i(w-z)$, so the centre of the
square on that side is $\frac{z+w}2-i\frac{w-z}2$. Hence $2p=(1+i)a+(1-i)b$, $2q=(1+i)b+(1-i)c$, $2r=(1+i)c+(1-i)d$, $2s=(1+i)d+(1-i)a$
for the squares on $AB$, $BC$, $CD$, $DA$. Using $i(1+i)=-(1-i)$ and $i(1-i)=1+i$,
\[
2i(s-q)=-(1-i)d+(1+i)a+(1-i)b-(1+i)c=-2(r-p).
\]
So $r-p=-i(s-q)$: the segments $PR$ and $QS$ have the same length and are perpendicular.
\end{solution}

\begin{solution}{ex:miquel}
\idea Define the point as the intersection of two of the circles and show that it is on the third, by adding angles.
\writeup Let the circles $AEF$ and $BFD$ meet at $F$ and at a second point $M$; assume $M$ is inside the triangle (the other cases are
the same with directed angles). In the cyclic quadrilateral $AFME$, $\angle FME=180^\circ-A$; in $BDMF$, $\angle FMD=180^\circ-B$. So
$\angle DME=360^\circ-(180^\circ-A)-(180^\circ-B)=A+B=180^\circ-C$, and the quadrilateral $CDME$ has opposite angles adding up to $180^\circ$:
it is cyclic, and $M$ lies on the circle $CDE$.
\end{solution}

\begin{solution}{ex:steinerlehmus}
\idea Equality of two expressions is hard to use directly; assume the triangle is not isosceles and show that one bisector is strictly
longer.
\writeup Let $a,b,c$ be the sides. The bisector from $B$ meets $AC$ at $D$ with $AD=\frac{bc}{a+c}$ and $DC=\frac{ab}{a+c}$ (the bisector
divides the opposite side in the ratio of the adjacent sides). Stewart's theorem, $a^2\,AD+c^2\,DC=b\,(BD^2+AD\cdot DC)$, which follows from the
cosine law in the triangles $ABD$ and $CBD$ (the cosines of the two angles at $D$ are opposite), gives
\[
t_b^2=BD^2=\frac{a^2c+c^2a}{a+c}-\frac{ab^2c}{(a+c)^2}=ac\Bigl(1-\frac{b^2}{(a+c)^2}\Bigr),
\]
and similarly $t_c^2=ab\bigl(1-\frac{c^2}{(a+b)^2}\bigr)$. If $t_b=t_c$, dividing by $a$ gives
\[
c-b=bc\Bigl(\frac b{(a+c)^2}-\frac c{(a+b)^2}\Bigr).
\]
If $c>b$, the left side is positive, but $b<c$ and $(a+c)^2>(a+b)^2$ make the right side negative. If $c<b$, the signs are reversed in
the same way. So $b=c$.
\end{solution}

\begin{solution}{ex:gausslucas}
\idea If $p'(w)=0$, the logarithmic derivative vanishes at $w$, and that is a balance of forces pulling $w$ towards the roots.
\writeup Write $p(z)=c\prod_k(z-z_k)$. If $p'(w)=0$ and $w$ is a root of $p$, there is nothing to prove. Otherwise
$0=\frac{p'(w)}{p(w)}=\sum_k\frac1{w-z_k}=\sum_k\frac{\overline{w-z_k}}{|w-z_k|^2}$. Conjugating and putting $m_k=|w-z_k|^{-2}>0$, we get
$\sum_km_k(w-z_k)=0$, so $w=\sum_km_kz_k/\sum_km_k$ is a convex combination of the roots. For a cubic, Marden's theorem says more: the two
roots of $p'$ are the foci of the ellipse inscribed in the triangle of the roots and tangent to its sides at their midpoints.
\end{solution}

\begin{solution}{ex:routh}
\idea Affine maps multiply all areas by the same factor and preserve ratios along lines, so we may choose the most convenient triangle.
\writeup Take $A=(0,0)$, $B=(1,0)$, $C=(0,1)$, of area $\frac12$, and $D=(\frac23,\frac13)$, $E=(0,\frac23)$, $F=(\frac13,0)$ on $BC$, $CA$, $AB$
with $BD:DC=CE:EA=AF:FB=1:2$. Solving pairs of linear equations, $AD\cap BE=(\frac47,\frac27)$, $BE\cap CF=(\frac17,\frac47)$ and
$CF\cap AD=(\frac27,\frac17)$. By Proposition~\ref{prop:tests}(3), the triangle they form has area
$\frac12\bigl|(-\frac37)(-\frac17)-\frac27(-\frac27)\bigr|=\frac1{14}$, one seventh of $\frac12$. Routh's theorem gives the general ratio: for
ratios $\frac{BD}{DC}=x$, $\frac{CE}{EA}=y$, $\frac{AF}{FB}=z$ it is $\frac{(xyz-1)^2}{(xy+x+1)(yz+y+1)(zx+z+1)}$, which is $\frac17$ for $x=y=z=\frac12$.
\end{solution}

\begin{solution}{ex:latticeequilateral}
\idea Areas of lattice triangles are rational, areas of equilateral triangles with lattice vertices would be $\sqrt3$ times a rational.
\writeup By Proposition~\ref{prop:tests}(3), the area is $\frac12|u_1v_2-u_2v_1|$ with integers $u_i,v_i$. An equilateral triangle of side $s$ has
area $\frac{\sqrt3}4s^2$, and $s^2$ is an integer (a sum of two squares). So $\frac{\sqrt3}4s^2$ would be rational with $s^2\ne0$, making $\sqrt3$
rational, which it is not.
\end{solution}

\begin{solution}{ex:cevaparallel}
\idea Parallel lines through the vertices cut the side lines proportionally, so the ratios can be read off one line.
\writeup Suppose $AD\parallel BE\parallel CF$, and choose coordinates in which these lines are vertical. Projecting a line onto the $x$-axis
along vertical lines multiplies all lengths on it by the same factor, so signed ratios on a non-vertical line can be read on the
$x$-coordinates. The side lines are not vertical (if $BC$ were vertical, the vertical line $AD$ would not meet it), and the
$x$-coordinates $x_A,x_B,x_C$ are distinct (if $x_A=x_B$, the cevian from $A$ would pass through $B$, so $D=B$). Since $D$ has
$x$-coordinate $x_A$, $E$ has $x_B$ and $F$ has $x_C$,
\[
\frac{BD}{DC}\cdot\frac{CE}{EA}\cdot\frac{AF}{FB}=\frac{x_A-x_B}{x_C-x_A}\cdot\frac{x_B-x_C}{x_A-x_B}\cdot\frac{x_C-x_A}{x_B-x_C}=1.
\]
\end{solution}

\section*{Sources and further reading}

\begin{itemize}
\item H.~S.~M.~Coxeter and S.~L.~Greitzer, \emph{Geometry Revisited}, Mathematical Association of America, 1967: the classical triangle
geometry of this chapter, with synthetic proofs; Section~3.3 tells the Napoleon story.
\item Proclus, \emph{A Commentary on the First Book of Euclid's Elements}, translated by G.~R.~Morrow, Princeton University Press, 1970
(the ``royal road'' of the epigraph).
\item J.~E.~Wetzel, \emph{Converses of Napoleon's theorem}, American Mathematical Monthly 99 (1992), 339--351 (including the history of the
theorem and its first appearance in \emph{The Ladies' Diary} for 1825).
\item J.~P.~Hogendijk, \emph{Al-Mu'taman ibn H\=ud, 11th century king of Saragossa and brilliant mathematician}, Historia Mathematica 22
(1995), 1--18.
\item L.-S.~Hahn, \emph{Complex Numbers and Geometry}, Mathematical Association of America, 1994; T.~Andreescu and D.~Andrica,
\emph{Complex Numbers from A to\dots Z}, 2nd ed., Birkh\"auser, 2014.
\item F.~Morley, \emph{On the intersections of the trisectors of the angles of a triangle}, Journal of the Mathematical Association of
Japan for Secondary Mathematics 6 (1924); C.~O.~Oakley and J.~C.~Baker, \emph{The Morley trisector theorem}, American Mathematical Monthly 85 (1978), 737--745 (history and
many proofs).
\item J.~H.~Conway, \emph{The power of mathematics}, in \emph{Power}, Cambridge University Press, 2006 (the backwards proof of Morley's
theorem); A.~Connes, \emph{A new proof of Morley's theorem}, Publications Math\'ematiques de l'IH\'ES 88 (1998), 43--46.
\item L.~Euler, \emph{Solutio facilis problematum quorundam geometricorum difficillimorum}, Novi Commentarii Academiae Scientiarum
Petropolitanae 11 (1767), 103--123 (the Euler line); K.~W.~Feuerbach, \emph{Eigenschaften einiger merkw\"urdigen Punkte des geradlinigen
Dreiecks}, N\"urnberg, 1822.
\item C.~Kimberling, \emph{Encyclopedia of Triangle Centers}, online, which lists tens of thousands of triangle centres with their
properties.
\item M.~Marden, \emph{Geometry of Polynomials}, American Mathematical Society, 1966 (Gauss--Lucas and Marden's theorem).
\end{itemize}
